\chapter*{Conclusion générale}
\addcontentsline{toc}{chapter}{Conclusion générale}

Ce rapport est parti d'un constat simple à énoncer mais difficile à concilier dans la pratique : la gestion du risque de crédit oblige une banque à arbitrer en permanence entre la maîtrise du risque de défaut, l'anticipation des pertes qu'il peut engendrer, la constitution d'un capital suffisant pour absorber les pertes inattendues, et la rentabilité des financements accordés. Le projet mené au sein de l'entité AI~\&~Data d'EY~Tunisie visait à construire une plateforme qui rapproche effectivement ces quatre préoccupations, plutôt que de les traiter séparément, au service de deux usages complémentaires, l'évaluation d'une nouvelle demande de crédit et le suivi d'un portefeuille déjà constitué.

Ce travail a mobilisé une chaîne complète, des données jusqu'à la décision. Une base de données de comptes de dépôt réels a d'abord servi de point d'ancrage à la construction d'historiques de crédit synthétiques, nettoyés et contrôlés, sur lesquels des modèles de probabilité de défaut et de perte en cas de défaut ont ensuite été développés et comparés selon une démarche champion/challenger, pour les segments Particuliers et Entreprises. L'exposition au moment du défaut, déterminée par des règles plutôt que par un modèle supplémentaire, a permis de calculer les pertes attendues selon IFRS~9, en distinguant l'horizon retenu pour le suivi de portefeuille de celui, simplifié, retenu à l'Origination. Le capital nécessaire pour absorber les pertes inattendues a été évalué selon deux approches délibérément distinctes, un forfait opérationnel pour chaque décision individuelle et une simulation Monte Carlo pour le pilotage du portefeuille, jamais combinées dans un même calcul de rentabilité. Plusieurs analyses complémentaires, segmentation du portefeuille, stress testing, projection de son évolution et simulation d'une politique d'octroi durcie, sont venues enrichir ce pilotage, avant qu'un mécanisme de décision explicite, fondé sur des seuils de probabilité de défaut, de RAROC et de stage définis dans des paramètres gouvernés, ne relie ces indicateurs à une recommandation explicable et traçable, tout en préservant l'intervention humaine dans les situations où une validation est requise. L'ensemble a été traduit en une architecture logicielle réelle, deux parcours temps réel et batch partageant des composants communs sans en confondre les usages, puis effectivement réalisé : une plateforme avec ses interfaces par rôle, sa gouvernance et sa traçabilité chaînée par hachage, et un assistant conversationnel qui interroge ces mêmes résultats sans jamais recalculer un indicateur de risque ni se substituer au mécanisme de décision. Cette réalisation a enfin été confrontée à une campagne de validation automatisée, portant sur 1\,345 exécutions de tests et de scénarios, dont les résultats et les limites sont discutés ci-dessous plutôt que résumés à un seul taux de réussite.

Ce travail comporte cependant des limites qu'il convient de rappeler sans les minimiser. Les historiques de crédit et de remboursement exploités tout au long de ce rapport restent des données synthétiques, construites à partir d'une base de comptes de dépôt réelle mais sans historique de défaut bancaire sous-jacent : les performances des modèles, aussi rigoureusement évaluées soient-elles sur ce périmètre, ne constituent donc pas une preuve de performance prédictive en conditions réelles, ni une validation indépendante au sens où l'exigerait un cadre réglementaire. Plusieurs simplifications méthodologiques ont par ailleurs été assumées plutôt que dissimulées : l'évaluation du modèle de perte en cas de défaut reste in-sample, faute d'un échantillon Entreprises suffisant pour une validation hors échantillon ; la perte attendue à l'Origination retient systématiquement un horizon de douze mois quel que soit le stage du client, par simplification de l'implémentation plutôt que par application intégrale de la norme ; et les deux bases de capital, forfaitaire et issue de la simulation Monte Carlo, ne sont à ce jour jamais réconciliées dans un même indicateur de rentabilité. La campagne de validation présentée au chapitre 6 a par ailleurs mis au jour trois anomalies réelles, non corrigées à ce stade : un pipeline d'extraction documentaire qui ne signale pas explicitement l'absence d'un fichier, des anomalies observées dans les scénarios de demande associés à un profil particulier de garantie, se traduisant selon le parcours d'exécution par des indicateurs financiers nuls ou une erreur de traitement, et une rupture ponctuelle de la chaîne de hachage du journal d'audit lors de deux écritures strictement simultanées. Enfin, cette plateforme n'a été ni déployée ni exploitée en environnement de production bancaire : sa réalisation démontre un fonctionnement logiciel réel, pas une mise en service opérationnelle.

Ces limites dessinent directement les perspectives d'évolution les plus pertinentes pour ce projet. La disponibilité d'historiques réels de défaut et de remboursement permettrait de reconduire l'ensemble de la démarche de modélisation sur des données dont la performance prédictive serait effectivement vérifiable, et d'envisager une validation hors échantillon et hors période pour le modèle de perte en cas de défaut, aujourd'hui limité par la taille de l'échantillon Entreprises. La réconciliation des deux bases de capital, aujourd'hui délibérément disjointes, et l'extension de la différenciation par stage au calcul de la perte attendue à l'Origination, constitueraient des enrichissements méthodologiques cohérents avec les limites déjà identifiées plutôt que des refontes. Les trois anomalies relevées lors de la campagne de validation gagneraient à être corrigées et rejouées, notamment celle touchant à l'intégrité du journal d'audit, dont la portée dépasse le seul cadre logiciel. Une validation en environnement bancaire réel, enfin, au-delà des tests logiciels présentés dans ce rapport, resterait le préalable nécessaire à toute évolution vers un usage opérationnel.

Au terme de ce travail, la plateforme développée constitue une démonstration cohérente et vérifiée de ce que peut apporter le rapprochement de l'évaluation du risque, de la quantification financière et de l'aide à la décision, sur le périmètre et avec les limites qui ont été précisément délimités tout au long de ce rapport. C'est cette honnêteté méthodologique, plus que l'exhaustivité des fonctionnalités réalisées, qui en constitue l'apport principal.
